\paragraph{\trConsistent{}地选取\trStriplikeEnd。}\label{sub:consistent-striplike-ends}
与前面一样，我们作\trStriplikeEnd{}的一个\trUniversal{}选取，并考虑与稳定的 $d$ \trLeafedTree{}$T$ 相关联的\trGluingConstruction。所得\trFamily{} \eqref{eq:gluing-family} 带有两组\trStriplikeEnd，它们\trApriori{}地看可能不同。第一组是通过\trGluingConstruction{}从每个 $\scr{S}^{|v|}$ 上的\trUniversal{}选取承袭而来；第二组则是经由\trClassifyingMap{} \eqref{eq:gluing-map} 从 $\scr{S}^{d+1}$ 上的\trUniversal{}选取\trPullback{}而来。如果对每棵 $T$，都存在一个包含 $\{0\}\times\prod_v\scr{R}^{|v|}$ 的开子集 $\bar{U}\subset\bar{\scr{R}}$，使得上述两组\trStriplikeEnd{}在 $U=\bar{U}\cap\scr{R}$ 上相同，那么我们就称\trStriplikeEnd{}的这个\trUniversal{}选取是\textbf{\trConsistent}（consistent）的。今后，我们把这样的 $U$ 或 $\bar{U}$ 称为\trGluingParameter{}充分小的子集。
\begin{lemma}\label{thm:consistent-ends-exist}
\trStriplikeEnd{}的\trConsistent{}\trUniversal{}选取是存在的。
\end{lemma}
这件事的证明简单却有些乏味，所以我们只解释其基本思想。第一个值得注意之点是，\trConsistency{}条件是归纳式的。设 $T$ 是一棵至少有两个\trVertex{}的 $d_0$ 叶稳定\trTree（一叶的\trLeafedTree{}给出空的\trConsistency{}条件）。$T$ 的\trConsistency{}借助诸 $\{\epsilon^{d+1}_k\}$（$d=|v|-1<d_0$）部分地确定了 $\scr{S}^{d_0+1}$ 上的诸\trStriplikeEnd{}$\{\epsilon^{d_0+1}_k\}$。下一个观察是，单个\trTree{}$T$ 的\trConsistency{}条件是可以满足的。细察\cref{thm:dm-smooth-structure}的证明即可看出，把 $\bar\gamma^T$ 限制到\trGluingParameter{}很小的子集 $\bar{U}$ 上，是一个\trOpenEmbedding，其像是 $\bar{\scr{R}}^{d_0+1}$\footnote{译注：原书此处印作 $\bar{\mathscr{R}}^{e+1}$；疑应为 $\bar{\mathscr{R}}^{d_0+1}$（$T$ 是 $d_0$ 叶\trTree{}，$\bar\gamma^T$ 的靶空间为 $\bar{\mathscr{R}}^{d_0+1}$，上标 $e$ 在原文中并无定义）；此处径改。——译者注}中 $T$ \trStratum{}的一个邻域。特别地，$\gamma^T|U$ 也是\trEmbedding。于是，若假设诸 $\{\epsilon^{d+1}_k\}$（$d<d_0$）已经取定，那么当然可以选取 $\{\epsilon^{d_0+1}_k\}$，使得它在 $\gamma^T|U$ 下的\trPullback——对单独一个具体的 $T$ 而言——与由\trGluing{}过程所得的 $\scr{S}|U$ 上那组\trStriplikeEnd{}相同。
论证的最后一步，是核查由各个 $T$ 所施加的条件并不互相矛盾。这又是一个递归论证，其基础是\trDeligneMumfordStasheff{}空间中\trBoundaryStratum{}的组合学。为具体起见，我们在 $d_0=4$ 这一最简单的非平凡情形下解释它。设我们已经选定了 $\{\epsilon^3_k\}$ 与 $\{\epsilon^4_k\}$，且其选取方式使得\trConsistency{}条件对稳定的 3 叶\trTree{}成立。带两个\trVertex{}的稳定 4 叶\trTree{}$T$ 共有五棵（其\trValency{}必然为 3 与 4），给出五个扩张\trGluingMap
\begin{equation*}
\bar\gamma^T:(-1;0]\times\scr{R}^3\times\scr{R}^4=(-1;0]\times\{\mathrm{point}\}\times(-1;1)\hookrightarrow\bar{\scr{R}}^5.
\end{equation*}
$\bar{\scr{R}}^5$ 是一个五边形。限制到\trGluingParameter{}很小的子集 $\bar{U}_T$ 上之后，每个\trGluingMap{}都是一个\trCovering{}某个\trBoundaryFace{}之邻域的\trEmbedding（见\cref{fig:pentagon-gluing}）。\trConsistency{}要求 $\{\epsilon^5_k\}$ 在这些邻域与内部 $\scr{R}^5$ 的交集上由 $\{\epsilon^3_k\}$、$\{\epsilon^4_k\}$ 确定。需要证明的是，这五个条件可以同时成立。任意两个邻域 $\bar\gamma^{T_1}(\bar{U}_{T_1})$、$\bar\gamma^{T_2}(\bar{U}_{T_2})$ 要么不交，要么在五边形的某个\trVertex{}附近重叠，后者对应于一棵带三个\trVertex{}的稳定 4 叶\trTree{}$T_3$。在前一种情形，诸\trConsistency{}条件显然互不干扰；在后一种情形，利用\trGluingMap{}$\bar\gamma^{T_3}$ 以及 $\{\epsilon^3_k\}$ 与 $\{\epsilon^4_k\}$ 之间的\trConsistency，可以验证由 $T_1$ 与 $T_2$ 的\trGluingConstruction{}所诱导的诸\trStriplikeEnd{}在\trGluingParameter{}取小值时相同，于是与 $\{\epsilon^5_k\}$ 相关联的\trConsistency{}条件便可同时得到满足。
\begin{figure}[htbp]
\centering
\begin{tikzpicture}[line cap=round,line join=round,scale=0.82]
\path (0,2) coordinate (V1)
      (-1.902,0.618) coordinate (V2)
      (-1.176,-1.618) coordinate (V3)
      (1.176,-1.618) coordinate (V4)
      (1.902,0.618) coordinate (V5);
\path (0,0) -- (V1) coordinate[pos=0.80] (W1);
\path (0,0) -- (V2) coordinate[pos=0.80] (W2);
\path (0,0) -- (V3) coordinate[pos=0.80] (W3);
\path (0,0) -- (V4) coordinate[pos=0.80] (W4);
\path (0,0) -- (V5) coordinate[pos=0.80] (W5);
\fill[gray!40] (V1)--(V2)--(W2)--(W1)--cycle;
\fill[gray!40] (V2)--(V3)--(W3)--(W2)--cycle;
\fill[gray!40] (V3)--(V4)--(W4)--(W3)--cycle;
\fill[gray!40] (V4)--(V5)--(W5)--(W4)--cycle;
\fill[gray!40] (V5)--(V1)--(W1)--(W5)--cycle;
\path (V2)--(V1) coordinate[pos=0.79] (A1);
\path (V5)--(V1) coordinate[pos=0.79] (B1);
\path (V3)--(V2) coordinate[pos=0.79] (A2);
\path (V1)--(V2) coordinate[pos=0.79] (B2);
\path (V4)--(V3) coordinate[pos=0.79] (A3);
\path (V2)--(V3) coordinate[pos=0.79] (B3);
\path (V5)--(V4) coordinate[pos=0.79] (A4);
\path (V3)--(V4) coordinate[pos=0.79] (B4);
\path (V1)--(V5) coordinate[pos=0.79] (A5);
\path (V4)--(V5) coordinate[pos=0.79] (B5);
\path (0,0) -- (A1) coordinate[pos=0.80] (A1p);
\path (0,0) -- (B1) coordinate[pos=0.80] (B1p);
\path (0,0) -- (A2) coordinate[pos=0.80] (A2p);
\path (0,0) -- (B2) coordinate[pos=0.80] (B2p);
\path (0,0) -- (A3) coordinate[pos=0.80] (A3p);
\path (0,0) -- (B3) coordinate[pos=0.80] (B3p);
\path (0,0) -- (A4) coordinate[pos=0.80] (A4p);
\path (0,0) -- (B4) coordinate[pos=0.80] (B4p);
\path (0,0) -- (A5) coordinate[pos=0.80] (A5p);
\path (0,0) -- (B5) coordinate[pos=0.80] (B5p);
\fill[gray!70] (A1)--(B1)--(B1p)--(A1p)--cycle;
\fill[gray!70] (A2)--(B2)--(B2p)--(A2p)--cycle;
\fill[gray!70] (A3)--(B3)--(B3p)--(A3p)--cycle;
\fill[gray!70] (A4)--(B4)--(B4p)--(A4p)--cycle;
\fill[gray!70] (A5)--(B5)--(B5p)--(A5p)--cycle;
\draw (V1)--(V2)--(V3)--(V4)--(V5)--cycle;
\node at (0,0.05) {$\bar{\scr{R}}^5$};
\foreach \px/\py/\rot in {-1.539/2.118/-54, -2.490/-0.809/18, 0/-2.618/90,
                          2.490/-0.809/162, 1.539/2.118/234}{
  \begin{scope}[shift={(\px,\py)},rotate=\rot]
    \draw (-1.1,0)--(-0.35,0);
    \draw (-0.35,0)--(-0.05,0.42) (-0.35,0)--(0.5,0) (-0.35,0)--(-0.05,-0.42);
    \draw (0.5,0)--(0.85,0.48) (0.5,0)--(0.98,0) (0.5,0)--(0.85,-0.48);
  \end{scope}}
\fill[gray!40] (-1.15,-4.15) rectangle (-0.75,-3.82);
\node[anchor=west] at (-0.6,-3.985) {\small\trGluingMap{}的像};
\fill[gray!70] (-1.15,-4.72) rectangle (-0.75,-4.39);
\node[anchor=west] at (-0.6,-4.555) {\small 重叠};
\end{tikzpicture}
\caption{}\label{fig:pentagon-gluing}
\end{figure}
